\section{Génération, conservation et réalisation}
\label{nuclear:introduccion}
L’objet de l’holographie modulaire triadique est un état engendré par des positions et des opérations qui conserve l’information de son évolution. L’évaluation numérique et l’incidence de bord sont des réalisations corrélées de cet antécédent. La contraction de cylindres compatibles détermine des valeurs ; le prolongement de leurs incidences maintient des relations qui distinguent les états générateurs. Cette priorité de la construction gouverne aussi la mécanique dimensionnelle : action, durée, étendue et masse apparaissent au moyen de lecteurs dont le domaine conserve l’orientation, la retenue et l’histoire des opérations.

Le présent article étudie une question structurelle précise : ce qui est conservé lorsqu’une lecture réduit l’état, où réside le complément et comment reconstruire à partir de lui aussi bien la valeur que la dynamique. Une réponse complète exige trois niveaux. D’abord sont construites les émissions et leur raffinement. Ensuite est démontrée une identité de transport sur l’état réalisé. Enfin est déterminé ce qu’observe un lecteur concret et quelles conséquences physiques produit sa composition avec les lecteurs d’action et d’échelle. Le développement suit cet ordre dans chaque application.

\subsection{Thèses et énoncés directeurs
}
La conservation évolutive de l’unité acquiert une expression opératorielle. Pour deux transports \(B,C\) entre espaces munis de formes \(b,b'\), avec

\[
B^{*}b'B=C^{*}b'C=b,
\]
la représentation collective des neuf positions produit

\begin{equation}
\mathcal U(B,C)=\frac19
\begin{pmatrix}8B+C&\sqrt8(C-B)\\ \sqrt8(C-B)&B+8C\end{pmatrix},
\qquad
\mathcal U(B,C)^*(b'\oplus b')\mathcal U(B,C)=b\oplus b.
\label{nuclear:eq:rectora}
\end{equation}
La première colonne contient la lecture et la mémoire engendrées par un état sans mémoire entrante. La seconde colonne décrit la contribution de la mémoire qui intervient à nouveau. L’égalité conserve une forme positive lorsque \(b\) est un produit scalaire, et conserve l’aire orientée lorsque \(b\) est une forme symplectique. Ses deux interprétations utilisent des structures explicitement différentes.


Cette loi permet de démontrer quatre conséquences reliées. La première est la reconstruction stable du registre complet, y compris son terminal. La deuxième récupère, à partir de la mémoire dodécaphasique, la direction qui détermine le projecteur dimensionnel. La troisième prolonge la conservation aux aires, aux volumes et aux degrés extérieurs supérieurs. La quatrième identifie, dans une réalisation quantique déclarée, la partie de la mémoire qui altère la pureté d’une lecture partielle. Chaque conséquence s’obtient par une opération mathématique spécifiée sur la précédente.

\begin{figure}[htbp]
\centering
\begin{tikzpicture}[every node/.style={font=\small,align=center},b/.style={draw,rounded corners,inner sep=6pt,text width=39mm},>=Stealth]
\node[b] (g) {Émission et état\\APP--TRIT--TPK};
\node[b,right=12mm of g] (u) {Transport complet\\et forme conservée};
\node[b,right=12mm of u] (l) {Lecture et mémoire\\avec reconstruction
};
\node[b,below=12mm of g] (i) {Incidence et\\raffinement};
\node[b,below=12mm of u] (a) {Action et\\réalisation dimensionnelle};
\node[b,below=12mm of l] (s) {Covariance, spectre\\et entropie réduite};
\draw[->] (g)--(u);\draw[->] (u)--(l);\draw[->] (g)--(i);
\draw[->] (u)--(a);\draw[->] (l)--(s);\draw[->] (i)--(a);\draw[->] (a)--(s);
\end{tikzpicture}
\caption{Dépendance des réalisations. Les flèches représentent des opérations démontrées dans le corps de l’article ; la lecture scalaire conserve un antécédent plus riche que sa propre valeur.
}
\end{figure}

\subsection{Mémoire complète et observation partielle}
Pour une suite de compressions \(T_n\) avec des compléments \(D_n\), on écrit
\(P_0=I\), \(P_{n+1}=T_nP_n\). L’égalité locale
\(T_n^*T_n+D_n^*D_n=I\) implique
\[
I=A_\infty+\sum_{n\ge0}P_n^*D_n^*D_nP_n,
\qquad
A_\infty=\underset{N\to\infty}{\operatorname{s-lim}}P_N^*P_N.
\]
La reconstruction utilise \(A_\infty^{1/2}v\) et tous les \(D_nP_nv\). La convergence forte du gramien terminal suffit à la conservation ; elle se démontre par monotonie positive, indépendamment de la convergence du vecteur terminal. Le registre contient des amplitudes et des corrélations, tandis que ses intensités fournissent le bilan de norme. Cette distinction détermine l’information récupérable.


Dans un registre fini concret, l’incidence des douze positions dans les hexades de Witt a pour Gram
\(36I+30\mathbf1\mathbf1^*\). Sa normalisation conserve exactement la partie centrée du registre. En incorporant la moyenne, on retrouve ses douze coordonnées. L’application dimensionnelle ultérieure utilise cette information transportée et conserve la dépendance à l’égard du registre ordonné \(K\).

\subsection{Une conséquence spectrale calculable}
Soit \(J^2=-I\), \(J^{\mathsf T}=-J\), et soit \(H\) symplectique dans une réalisation canonique de dimension \(2d\). On définit
\[
D=\frac{\sqrt8}{9}(H-I),\qquad
D_{\mathrm{anti}}=\frac{D+JDJ}{2},\qquad
W=\frac{8I+HH^{\mathsf T}}9.
\]
L’article démontre

\begin{equation}
\boxed{-(JW)^2=I+[D,J][D,J]^{\mathsf T}
=I+4D_{\mathrm{anti}}D_{\mathrm{anti}}^{\mathsf T}.}
\label{nuclear:eq:espectro}
\end{equation}
Pour deux entrées gaussiennes pures identiques et le transport \(\mathcal U(I,H)\), la matrice \(W\) est la covariance réduite normalisée. Par conséquent,
\[
\operatorname{Tr}\rho_{\mathrm{vis}}^2
=\det\bigl(I+4D_{\mathrm{anti}}D_{\mathrm{anti}}^{\mathsf T}\bigr)^{-1/4}.
\]
La formule relie un opérateur de mémoire à la fonctionnelle de pureté de l’état réduit. La pureté globale reste égale à un. En particulier, le mélange dépend de la composante antilinéaire par rapport à la structure complexe transportée, et non d’une identification indifférenciée entre mémoire et entropie.


Le choix de la représentation gaussienne, de l’état d’entrée et du sous-système lu fait partie de l’énoncé. La généalogie fixe les transports qui sont composés ; le protocole fixe l’observation dont la distribution est calculée. Cette spécification permet une comparaison reproductible entre réalisations et sépare le théorème de toute attribution expérimentale ultérieure.

\subsection{Ordre de l’exposé}
Les chapitres initiaux développent intégralement les primitives communes, le prolongement, les régions numériques, le registre et l’incidence exceptionnelle. Suivent les lecteurs d’action et les représentations dimensionnelles. La deuxième partie réunit les preuves de conservation, de récupération, d’holonomie et d’extension réversible. La partie finale compose ces résultats avec des actions physiques et obtient le spectre de la mémoire réduite. La bibliographie situe les outils de représentation après la construction qu’ils réalisent.

La symétrie finie s’exprime au moyen d’actions et d’opérateurs d’entrelacement. La symétrie variationnelle continue s’exprime au moyen d’une action, d’une variation et de sa charge conservée. Le théorème de Noether s’applique dans ce second domaine ; le transport d’incidence exceptionnelle conserve son propre contenu algébrique. La relation entre les deux s’établit par les observables et les formes transportées, en maintenant explicites leurs différences de type.


% BEGIN R27_FRAMING_INTRO_ES
La conservation est aussi examinée avant l’introduction de représentations
linéaires. Le codage décimal de l’émetteur possède une congruence de
contrôle et un alphabet complet de \(42\) blocs, démontré au moyen
de ses signatures et de ses germes. Les incidences entre moitiés d’émissions unissent
les régions en conservant la provenance de chaque arête. La comparaison
de deux registres ayant un histogramme commun et une corrélation différente établit
une distinction concrète entre distribution et ordre. Les lois d’archivage
reçoivent cet ordre avec les amplitudes et les données terminales
qui figurent dans leurs énoncés.

La normalisation du logarithme multiplicatif est distinguée de son
inversion au moyen du critère du plus petit générateur positif. Sur la
clôture entière, les contrôles du sens de la retenue, de l’extrémité terminale,
de la rotation et de la perturbation locale maintiennent les autres données fixes.
La sélection de frontière applique successivement la saturation, la neutralité
duale et l’orientation. Ces opérations précèdent les lectures de
récupération et conservent leur domaine, de sorte que chaque inverse agit
sur les données effectivement produites par sa construction.
% END R27_FRAMING_INTRO_ES

% BEGIN R28_FRAMING_INTRO_ES
Le lien entre les registres est explicité au moyen du multigraphe de
l’émetteur et de son quotient de phase : cinq fermetures minimales réunissent les quatre
incidences requises pour chaque région. La dérivation de l’occupation
de frontière et la parité des lecteurs sous inversion orientée
complètent la description des données que reçoit le transport bilatéral
et qui maintiennent ses lectures de couplage.

% END R28_FRAMING_INTRO_ES
